\section{Method}
\label{sec:method}

\subsection{Problem Setting and Overview}

Let $\mathcal D_s=\{(x_j^s,y_j^s)\}_{j=1}^{n_s}$ denote labeled reference
data and $\mathcal X_t=\{x_i^t\}_{i=1}^{n_t}$ available unlabeled target queries, with
$x_j^s,x_i^t\in\mathbb R^d$. We consider classification with
$\mathcal Y=\{1,\ldots,L\}$ and regression with $\mathcal Y\subseteq\mathbb R$.
A frozen tabular foundation model
$f_\theta$ predicts each query from a labeled context constructed from
source examples and, for registered selection-shift constructors,
pseudo-labeled target examples. Ground-truth target labels are used only for
evaluation, and $f_\theta$ remains frozen at test time.
Reference and target examples follow $P_s(X,Y)$ and $P_t(X,Y)$, respectively.
We study both controlled changes to the accessible context and distribution
shifts with $P_t\neq P_s$, using $\mathcal D_s$, the current query, and any
target features required by the registered constructor at inference.

As illustrated in Figure~\ref{fig:tabace_overview}, \method{} uses one Base
branch and multiple alternative context branches while keeping $f_\theta$
frozen. Confidence and cross-context agreement determine the relative influence
of the alternatives, while the task-specific gate and aggregate context support
set the update magnitude relative to Base. The update therefore follows the
available query-level evidence.

\subsection{Configurable Context and Output Interface}
\label{sec:context_interface}

Let $\phi_0$ denote the Base configuration and
$\Phi=\{\phi_1,\ldots,\phi_R\}$ the alternative branch configurations chosen
from source-side information and application knowledge before target inference.
We write $\mathcal I(\Phi)=\{1,\ldots,R\}$ for their indices. Let
$F_{\phi_r}$ denote the complete prediction mapping induced by configuration
$\phi_r$. For an available query $x_i^t$,
\begin{equation}
  o_i^{(r)}
  =F_{\phi_r}(\mathcal D_s,\mathcal X_t,x_i^t;f_\theta),
  \qquad r=0,\ldots,R.
  \label{eq:context_prediction}
\end{equation}
The $\mathcal X_t$ argument supplies any target subset required by $\phi_r$ and
is ignored by source-only constructors.
For a TFM branch, this mapping composes a configurable context constructor
with the frozen predictor. A shared-context constructor forms
$C^{(r)}=G_{\phi_r}(\mathcal D_s,\mathcal X_t)$ once and reuses it across
queries, so that
$F_{\phi_r}(\mathcal D_s,\mathcal X_t,x_i^t;f_\theta)
=f_\theta(C^{(r)},x_i^t)$. A query-dependent construction forms
$C_i^{(r)}=G_{\phi_r}(\mathcal D_s,\mathcal X_t,x_i^t)$. A compatible conventional predictor uses the
constructed context as its branch-specific training data and exposes the same
output interface. For classification,
$o_i^{(r)}=p_i^{(r)}\in\Delta^{L-1}$ is a class-probability vector; for point regression,
$o_i^{(r)}=\hat y_i^{(r)}$ is scalar.

We write $o_i^{(0)}$ for the Base output.
Its constructor $G_{\phi_0}$ follows the backbone's standard sampling rule to
form the canonical context from the available labeled reference data, with its
selected features and labels unchanged. For
$r\geq1$, $G_{\phi_r}$ constructs a complementary context by selecting or
reweighting reference examples, selecting or transforming features, or
applying a prescribed source-derived label adjustment. A selection-shift
constructor may also attach Base pseudo-labels to target peers. For query
$x_i^t$, these rows satisfy
$\widehat{\mathcal D}_{t,-i}^{(0)}\subseteq
\{(x_j^t,\widetilde y_j^{(0)}):j\neq i\}$, where
$\widetilde y_j^{(0)}=\arg\max_{\ell\in\mathcal Y}p_{j,\ell}^{(0)}$. The
current query is excluded. Constructors use reference information and, when
registered, statistics of the available target subset. Target
labels remain unavailable to all constructors. When a constructor selects a
feature subset, the same
projection is applied to the context and the query. The resulting
configurations produce the alternative output bank
\begin{equation}
  \mathcal O_i(\Phi)=\{o_i^{(r)}:r\in\mathcal I(\Phi)\}.
  \label{eq:output_bank}
\end{equation}
Here, $\Phi$ specifies the context-construction rules and
$\mathcal O_i(\Phi)$ collects their query-dependent outputs.

\subsection{Context Aggregation by Confidence and Agreement}
\label{sec:evidence_aggregation}

For query $x_i^t$, the $R$ alternative contexts yield Eq.~\eqref{eq:output_bank}.
We derive task-compatible support scores and aggregate the predictions with a
shared normalized weighting rule.

For classification, let $p_i^{(r)}\in\Delta^{L-1}$ be the prediction from
context $r$. With source-learned $\eta_{Q,r},\eta_{S,r}\geq0$, we compute
\begin{equation}
  \begin{aligned}
    \mu_i&=\frac{1}{R}\sum_{s=1}^{R}p_i^{(s)},
    &Q_i^{(r)}&=1-\frac{H(p_i^{(r)})}{\log L},\\
    S_i^{(r)}&=1-\frac12\lVert p_i^{(r)}-\mu_i\rVert_1,
    &\rho_i^{(r)}&=\eta_{Q,r}Q_i^{(r)}+\eta_{S,r}S_i^{(r)}.
  \end{aligned}
  \label{eq:classification_evidence}
\end{equation}
Here $H(p)=-\sum_{\ell=1}^{L}p_\ell\log p_\ell$.
$Q_i^{(r)}$ measures normalized-entropy confidence and $S_i^{(r)}$ measures
agreement with the mean alternative prediction; $\rho_i^{(r)}$ increases when
both are high.

For regression, let $\hat y_i^{(r)}\in\mathbb R$ be the prediction from
context $r$. We define
\begin{equation}
  \begin{aligned}
    m_i&=\operatorname{median}\{\hat y_i^{(s)}\}_{s=1}^{R},
    &S_i^{(r)}&=\exp\!\left(-\frac{|\hat y_i^{(r)}-m_i|}{\tau_y}\right),\\
    \rho_i^{(r)}&=\eta_{S,r}^{\mathrm{reg}}S_i^{(r)},
    \qquad \eta_{S,r}^{\mathrm{reg}}\geq0.
  \end{aligned}
  \label{eq:regression_evidence}
\end{equation}
The median $m_i$ is a robust center and $\tau_y>0$ is determined from source
data; predictions far from this center receive less support.

Both tasks then use the same normalized weighting and aggregation rule:
\begin{equation}
  \begin{aligned}
    w_i^{(r)}&=\frac{\pi_r e^{\rho_i^{(r)}}}
      {\sum_{s=1}^{R}\pi_s e^{\rho_i^{(s)}}},\\
    \bar o_i&=\sum_{r=1}^{R}w_i^{(r)}o_i^{(r)},
    \qquad \sum_{r=1}^{R}w_i^{(r)}=1.
  \end{aligned}
  \label{eq:task_aware_aggregation}
\end{equation}
Here $\pi_r>0$ is the source-learned prior reliability of context $r$, with
$\sum_r\pi_r=1$. Thus $(o_i^{(r)},\bar o_i)=(p_i^{(r)},\bar p_i)$ for
classification and $(o_i^{(r)},\bar o_i)=(\hat y_i^{(r)},\bar y_i)$ for
regression. The aggregate $\bar o_i$ is then integrated with Base.

\subsection{Query-Dependent Integration with Base}
\label{sec:base_integration}

We integrate $\bar o_i$ with Base using a task-specific gate $U_i$ and aggregate
context support $E_i$. Here $U_i$ is intrinsic Base uncertainty for classification
and the Base-aggregate discrepancy for regression, while $E_i$ measures
alternative support:
\begin{equation}
  \begin{aligned}
    U_i^{\mathrm{cls}}&=\frac{H(p_i^{(0)})}{\log L},
    &E_i^{\mathrm{cls}}&=\sum_{r=1}^{R}w_i^{(r)}Q_i^{(r)}S_i^{(r)},\\
    U_i^{\mathrm{reg}}&=1-\exp\!\left(-\frac{|\bar y_i-\hat y_i^{(0)}|}{\tau_y}\right),
    &E_i^{\mathrm{reg}}&=\sum_{r=1}^{R}w_i^{(r)}S_i^{(r)}.
  \end{aligned}
  \label{eq:context_integration_evidence}
\end{equation}
Below, $U_i$ and $E_i$ denote the pair for the current task.

Both tasks use the same integration rule:
\begin{equation}
  \alpha_i=\alpha_{\max}U_iE_i,
  \qquad o_i^{\mathrm{final}}
  =(1-\alpha_i)o_i^{(0)}+\alpha_i\bar o_i.
  \label{eq:base_anchored_update}
\end{equation}
Since $U_i,E_i,\alpha_{\max}\in[0,1]$, the update remains close to Base whenever
the task-specific gate or context support is weak.

\paragraph{Outer-fold policy learning.}
Calibration episodes from outer-training datasets fit the context priors, score
coefficients, output scale, and integration cap. Each episode pairs a labeled
source context with a pseudo-target feature set whose labels enter only the
calibration loss. The fitted policy and $\Phi$ are fixed before held-out inference.
